\chapter{电平、开关、逻辑门与组合电路}

本章将“电和器件”“开关到门”“组合逻辑”合并为一条完整链路：连续变化的电压先被解释成稳定的 0/1；二极管和受控开关提供可预测的电流路径；开关网络形成 NOT、NAND、NOR、XOR 等门；门再组合成判断器、加法器、选择器和译码器。这样的组织方式有助于读者在同一章内建立从物理电平到组合电路的连续推导。

本章不展开完整的半导体物理，也不把内容处理成布尔代数速查表。它关注一个核心问题：一根真实导线上的电压，怎样逐步成为 CPU 可使用的 bit 计算。后续所有寄存器、ALU、数据通路、控制器、内存和 I/O，最终都必须满足这里建立的底层合同。

\begin{keyidea}
数字硬件的第一层抽象不是门，而是带余量的电平；数字硬件的第一层实现不是公式，而是受控电流路径。
\end{keyidea}

\section{一、从一根线得到可靠 bit}

电路里说“一根线是 1”，完整含义不是“线上写着数字 1”，而是“这个节点相对于参考地的电压落在接收端承认的高电平区间内”。节点可以理解成多条金属连接汇合成的同一个电气位置；参考地是整块电路共同约定的 0V 参考点；电压是两个点之间的差值。数字电路里说某根信号线是 0V、3.3V 或 5V，实际都是在说它相对于地的电压。

电流说明电荷正在沿通路流动。若输出节点被一条通路接向电源，它可能被拉到高电平；若被一条通路接向地，它可能被拉到低电平。若两条通路都断开，节点不会自动变成 0，而是可能暂时停在旧电压附近，再被漏电流、干扰、后级输入或测量探针慢慢推走。若上拉和下拉两条强通路同时打开，电源到地之间形成冲突电流，电平不可靠，还会增加功耗和发热。

需要先区分导线和程序变量。变量没有被赋值时，语言可以规定默认值；导线没有被驱动时，只剩物理过程。它可能因为寄生电容保存旧电压，表现为保留上一次状态；也可能被旁边切换的信号线、人体接近、探针输入阻抗或后级漏电流轻微影响。调试数字硬件时，不能把这种偶然稳定误认为设计正确。

以按钮输入为例。假设按钮按下表示“请求发生”。一种不完整的接法是按钮一端接信号线，另一端接电源：

\begin{lstlisting}
button pressed  -> signal connected to VCC
button released -> signal connected to nothing
\end{lstlisting}

按下时，信号线被电源路径拉高，该部分行为符合预期。松开时，信号线并没有被明确拉低，它只是断开了。这个节点可能保留先前电荷，也可能被旁边导线耦合出一定电压，还可能被后级输入漏电流拖向某个不确定位置。于是接收端可能读到 1、0 或过渡电压。问题的根源不是程序错误，而是电路没有为节点规定明确归宿。

正确的第一层想法是：任何要被后级判断的输出节点，都必须在每一种输入情况下被明确地拉向高电平或低电平。按钮松开时需要一条弱下拉路径把信号线拉向地；按钮按下时，按钮路径把它拉向电源。弱下拉不是为了在按钮按下时和按钮争夺电源，而是为了在按钮松开时给节点一个默认归宿。它太弱，节点下降慢、抗干扰差；它太强，按钮按下时浪费电流，甚至让高电平不够高。

\begin{longtable}{p{0.18\linewidth}p{0.31\linewidth}p{0.31\linewidth}p{0.12\linewidth}}
\toprule
场景 & 高电平路径 & 低电平路径 & 结论 \\
\midrule
按下 & 按钮把节点接向电源 & 弱下拉仍存在 & 可靠高 \\
松开 & 高电平路径断开 & 弱下拉把节点接向地 & 可靠低 \\
接触抖动 & 高电平路径快速断续 & 弱下拉持续存在 & 需要消抖 \\
导线断裂 & 按钮路径永远断开 & 若下拉还在，则为低 & 有默认值 \\
没有下拉 & 高电平路径断开 & 低电平路径也断开 & 悬空 \\
\bottomrule
\end{longtable}

该表给出硬件读法的雏形。它不只问逻辑上应该是 0 还是 1，还问路径在哪里、谁驱动谁、异常时会落到哪个默认状态。后面讲复位、总线仲裁和设备中断时，同样要用这种检查法：每个状态都要有明确来源，每个默认值都要有物理路径支持。

电平也不能理解成两个数学点。假设一块小电路用 5V 供电。线上的真实电压可能是 0V、0.17V、1.3V、2.6V、4.8V，也可能在切换过程中从 0V 慢慢升到 5V。硬件不能要求每一级都精确读出这个连续数值。更可靠的做法是规定两个稳定区间：低到足够接近地电位时解释为 0，高到足够接近电源电位时解释为 1，中间一段不作为长期稳定状态。

\begin{longtable}{p{0.24\linewidth}p{0.28\linewidth}p{0.38\linewidth}}
\toprule
电压范围 & 逻辑解释 & 设计意义 \\
\midrule
接近地电位 & 逻辑 0 & 小噪声不会轻易把它推成 1 \\
中间区域 & 不作为稳定状态 & 切换可以经过，但不能长期停留 \\
接近电源电位 & 逻辑 1 & 小噪声不会轻易把它拉成 0 \\
\bottomrule
\end{longtable}

可以用四个电压指标描述两级电路之间的合同。$V_\mathrm{OH}$ 是输出端保证的高电平下限，$V_\mathrm{OL}$ 是输出端保证的低电平上限；$V_\mathrm{IH}$ 是输入端承认为高的最低电压，$V_\mathrm{IL}$ 是输入端承认为低的最高电压。若一个电路保证 $V_\mathrm{OH}$ 高于下一级需要的 $V_\mathrm{IH}$，中间差额就是高电平噪声余量；若 $V_\mathrm{OL}$ 低于下一级允许的 $V_\mathrm{IL}$，中间差额就是低电平噪声余量。

\begin{longtable}{p{0.18\linewidth}p{0.32\linewidth}p{0.4\linewidth}}
\toprule
符号 & 含义 & 读法 \\
\midrule
$V_\mathrm{OH}$ & 输出保证为高时的最差电压 & 前一级说：我给出的 1 至少高到这里 \\
$V_\mathrm{IH}$ & 输入承认为高的最低电压 & 后一级说：至少给到这里，我才认作 1 \\
$V_\mathrm{OL}$ & 输出保证为低时的最差电压 & 前一级说：我给出的 0 至多高到这里 \\
$V_\mathrm{IL}$ & 输入承认为低的最高电压 & 后一级说：低到这里以下，我才认作 0 \\
\bottomrule
\end{longtable}

数字抽象并不是忽略模拟世界，而是把模拟世界约束在合同之内。只要所有连接都满足这些最差条件，后续讨论 bit、真值表和指令编码才有可靠基础。若这些条件不满足，系统层面可能出现间歇性现象：同一块板有时能启动、有时不能；按键有时触发两次；某条总线在高温下出错；更换更长排线后原本稳定的电路开始误判。这些现象不一定来自逻辑表达式错误，也可能来自电平合同被破坏。

输出还必须能驱动下一级。一个节点如果只能在空载时看起来像 1，一接上后级就被拖低，它不是可靠数字输出。后级输入、板级走线、封装和探针都有寄生电容。把一个节点从 0 拉到 1，本质上是在给这个电容充电；把它从 1 拉到 0，本质上是在放电。驱动路径越弱，等效电阻越大；负载越多，等效电容越大；上升和下降越慢。

\begin{lstlisting}
larger load capacitance -> slower edge
weaker driver           -> slower edge
slower edge             -> less time margin before sampling
\end{lstlisting}

所以示波器上看到的真实信号不是理想直角。电压边沿有斜率，可能有过冲、振铃、平台和噪声。逻辑 0/1 的背后是一条随时间变化的电压曲线，接收端只是在某个时刻把它归类为 0 或 1。若这个时刻到来之前电压还没有进入可靠区间，最终会稳定也不够。

\section{二、从器件到受控电流路径}

二极管可以先理解成方向性很强的器件。在合适偏置下，它比较容易让电流沿一个方向通过；反向偏置时，它基本阻止电流通过。这个模型不需要一开始展开半导体物理，却已经足够说明硬件设计的一条基本路线：利用器件物理特性，制造可预测的电流路径。

二极管的“单向”不是理想数学开关。正向导通通常需要一定压差，导通后还有压降；反向并非绝对没有电流，而是漏电很小；电压过高时还可能击穿。初读数字硬件时可以先用理想模型建立方向感，但不能忘记真实器件会吃掉电压余量、引入延迟和功耗。

\begin{longtable}{p{0.2\linewidth}p{0.34\linewidth}p{0.36\linewidth}}
\toprule
二极管事实 & 对数字电路的影响 & 读书时要追问 \\
\midrule
有方向性 & 可以限制电流路径 & 这条路径在何种输入下导通 \\
有正向压降 & 输出电平会损失一部分余量 & 级联后高电平还够不够高 \\
有漏电和极限 & 不能把反向阻断看成无限好 & 极端温度、电压下是否仍可靠 \\
不能主动放大 & 不能恢复变弱的电平 & 下一级是否需要缓冲或开关恢复 \\
\bottomrule
\end{longtable}

二极管逻辑可以表达某些简单关系。若有两路输入想共同点亮一个指示节点，任意一路为高时都能通过二极管把输出节点拉高，输出就表现出 OR 的味道：

\begin{lstlisting}
if A can drive through a diode:
  Y may be pulled high
if B can drive through a diode:
  Y may be pulled high
\end{lstlisting}

该例能够说明方向性逻辑，但也暴露出局限。第一，二极管会带来压降，输出高电平可能比输入高电平低一截；级联多了，电平余量会越来越紧。第二，二极管只能提供方向性，不能主动把一个偏弱的高电平重新恢复成强高电平。第三，它不方便实现取反，而取反是构造通用逻辑必须具备的能力。

取反为什么重要？因为许多条件天然需要反向解释。传感器未触发时允许运行，触发时禁止运行；按钮未按下时保持，按下时改变。若电路只能表达“有一路导通”，却不能把 1 变成 0、把 0 变成 1，那么可组合的逻辑很快就受限。要继续向前，需要一个信号能控制另一条电流路径是否导通。

晶体管的细节很深，但作为第一层硬件模型，可以先把它当成受控开关。控制端不是直接给输出供电，而是决定输出节点到某条电源路径之间是否接通。受控开关和普通机械开关最大的差别，是控制信号本身也是电信号。按钮需要人按，晶体管可以由前一级逻辑输出控制。于是硬件有了“信号控制信号”的递归能力：一个门的输出可以控制下一批开关，下一批开关再形成新的输出。

反相器可以作为第一种完整逻辑门来理解。输入低时，上拉路径导通、下拉路径断开，输出被拉高；输入高时，上拉路径断开、下拉路径导通，输出被拉低。注意这里不是“输入低所以输出自动高”，而是“输入低使某个高电平路径导通”。把每一句逻辑描述翻译成路径描述，能避免很多硬件误解。

\begin{longtable}{p{0.18\linewidth}p{0.29\linewidth}p{0.29\linewidth}p{0.14\linewidth}}
\toprule
输入 & 上拉路径 & 下拉路径 & 输出 \\
\midrule
低电平 & 导通 & 断开 & 被拉高 \\
高电平 & 断开 & 导通 & 被拉低 \\
过渡中 & 可能部分导通 & 可能部分导通 & 暂时不应采样 \\
\bottomrule
\end{longtable}

第三行很重要。真实输入从低到高不是瞬间跳变，在过渡区间里，上拉和下拉可能都不处于理想开/关状态。短时间内出现电源到地的电流并不奇怪，这就是动态功耗和短路功耗的一部分。只要过渡足够快、采样避开中间区域，数字抽象仍然成立。若输入边沿太慢，电路可能在中间区域停留过久，功耗和误判风险都会上升。

把开关串联，可以表达“两个条件都满足才有通路”；把开关并联，可以表达“任意一个条件满足就有通路”。以下拉网络为例，两个受控开关串联，只有 A 和 B 都让开关导通时，输出节点才会被拉向低电平；两个开关并联，只要 A 或 B 有一个导通，输出节点就会被拉低。上拉网络通常要和下拉网络互补，保证输出在稳定输入下要么被拉高，要么被拉低，而不是两边都断或两边都强通。这种“互补路径”思想，后面会自然走到 CMOS 逻辑。

\section{三、从路径表得到逻辑门}

门电路的第一步不是画符号，而是把输入组合列完整。两个输入 A、B 只有四种组合。电路必须对每一种组合给出明确输出，而且每一行都要满足前面的要求：输出节点不能悬空，不能上下冲突，必须能在规定时间内进入可靠电平区间。

以安全允许灯为例。A 表示保护盖合上，B 表示急停未按下。允许灯亮的条件很严格：任一条件不满足，输出都必须为 0。

\begin{longtable}{p{0.12\linewidth}p{0.18\linewidth}p{0.16\linewidth}p{0.42\linewidth}}
\toprule
A & B & 允许灯 & 为什么 \\
\midrule
0 & 0 & 0 & 盖子没合上，急停也不是可运行状态 \\
0 & 1 & 0 & 急停条件满足，但盖子没合上 \\
1 & 0 & 0 & 盖子合上了，但急停条件不满足 \\
1 & 1 & 1 & 两个安全条件同时满足 \\
\bottomrule
\end{longtable}

该表对应 AND 的行为。它不是抽象数学定义，而是“两个条件都满足才允许”的硬件约束。若某个实现让 A=1、B=0 时灯也亮，那就不是 AND 近似得不够好，而是安全行为错了。真值表的价值在于防止“只凭自然语言写电路”。“两个条件都满足”看起来简单，但实现时很容易漏掉某一行。

再看故障灯。A 表示温度报警，B 表示电流报警。故障灯的规则反过来：只要有一个坏消息出现，就要亮。

\begin{longtable}{p{0.12\linewidth}p{0.18\linewidth}p{0.16\linewidth}p{0.42\linewidth}}
\toprule
A & B & 故障灯 & 为什么 \\
\midrule
0 & 0 & 0 & 没有任何故障证据 \\
0 & 1 & 1 & 电流报警已经足够触发故障 \\
1 & 0 & 1 & 温度报警已经足够触发故障 \\
1 & 1 & 1 & 两个故障同时存在，仍应报警 \\
\bottomrule
\end{longtable}

该表对应 OR 的行为。AND 和 OR 的差别不在符号长什么样，而在需求不同：一个要求所有条件同时成立，一个要求任意证据成立。NOT 只有一个输入，但同样要列完整。若 A 表示“禁止信号”，输出 Y 表示“允许动作”，那么 A=0 时 Y=1，A=1 时 Y=0。它解决的是反向解释：输入出现时，输出反而关闭。

\begin{warningbox}
不要把“0 是假、1 是真”机械套到每个硬件信号上。某些信号是低有效，例如 \code{reset_n}、\code{chip_select_n}，低电平才表示动作发生。读这类信号时，先写真值表，再谈门电路。
\end{warningbox}

真值表说清楚“应该怎样”，路径表说明“为什么会这样”。用受控开关看 AND，最自然的结构是串联。两个开关串在同一条路径上，只有 A 和 B 都让自己的开关闭合时，整条路径才通。OR 更像并联。A 控制一条路径，B 控制另一条路径。只要任意一路闭合，输出就能被拉向目标电平。

\begin{longtable}{p{0.12\linewidth}p{0.12\linewidth}p{0.34\linewidth}p{0.32\linewidth}}
\toprule
A & B & 串联路径状态 & AND 输出为什么成立或不成立 \\
\midrule
0 & 0 & 两个开关都断开 & 没有完整导通路径 \\
0 & 1 & A 断开，B 闭合 & 串联路径仍被 A 切断 \\
1 & 0 & A 闭合，B 断开 & 串联路径仍被 B 切断 \\
1 & 1 & 两个开关都闭合 & 从源到输出的路径完整 \\
\bottomrule
\end{longtable}

如果只画一条“让输出为 1 的路径”，还没有完成门电路设计。输出节点不能靠悬空表示 0。因此每一行真值表还要补一张路径表，说明高电平路径和低电平路径的状态。一个合法门在稳定输入下，不应该让输出悬空，也不应该让上拉和下拉强路径同时导通。

从开关网络看，NAND 和 NOR 往往比 AND 和 OR 更像基础积木。原因是它们天然带反相输出，更容易形成互补的上拉和下拉结构：当下拉网络导通时输出被拉低；当下拉网络不导通时，上拉网络负责把输出拉高。

\begin{longtable}{p{0.1\linewidth}p{0.1\linewidth}p{0.22\linewidth}p{0.46\linewidth}}
\toprule
A & B & NAND 输出 & 开关直觉 \\
\midrule
0 & 0 & 1 & 下拉串联路径断开，上拉负责拉高 \\
0 & 1 & 1 & A 断开，下拉路径不完整 \\
1 & 0 & 1 & B 断开，下拉路径不完整 \\
1 & 1 & 0 & A、B 都闭合，下拉路径完整，输出被拉低 \\
\bottomrule
\end{longtable}

NOR 可以用同样方式读。NOR 的输出只有在 A=0 且 B=0 时为 1；只要 A 或 B 为 1，输出就为 0。若以下拉网络看，A、B 并联，只要任一路输入为 1，就能把输出拉低；只有两路都不导通时，上拉网络把输出拉高。

只用 NAND 就能构造 NOT、AND、OR：

\begin{lstlisting}
NOT A = A NAND A
A AND B = (A NAND B) NAND (A NAND B)
A OR B = (A NAND A) NAND (B NAND B)
\end{lstlisting}

这说明 NAND 在功能上完备，不说明所有实现都应该只用 NAND。真实设计会在功能正确的前提下，考虑门级数量、路径延迟、面积、功耗和负载。逻辑式等价，只说明稳定输出行为相同；它不保证延迟、功耗和驱动能力相同。

XOR 的输出在两个输入不同的时候为 1，相同的时候为 0。它解决的是“检测两个 bit 是否不同”这个具体问题。

\begin{longtable}{p{0.16\linewidth}p{0.16\linewidth}p{0.24\linewidth}p{0.34\linewidth}}
\toprule
A & B & A XOR B & 解释 \\
\midrule
0 & 0 & 0 & 两个输入相同 \\
0 & 1 & 1 & 两个输入不同 \\
1 & 0 & 1 & 两个输入不同 \\
1 & 1 & 0 & 两个输入相同 \\
\bottomrule
\end{longtable}

可以把 XOR 拆成两种会输出 1 的情况：A=1 且 B=0，或者 A=0 且 B=1：

\begin{lstlisting}
A XOR B = (A AND NOT B) OR (NOT A AND B)
\end{lstlisting}

这条式子不要只当代数记忆。它对应两条证据路径：第一条证明“只有 A 为 1”，第二条证明“只有 B 为 1”。最后的 OR 合并这两条证据。若两个都为 0，没有证据成立；若两个都为 1，两条“只有一个为 1”的证据也都不成立。XOR 后面会反复出现：两个 bit 相加时，和位就是“两个输入是否不同”；比较两个 bit 时，XOR 可以告诉我们它们是否不同；接一个 NOT，就能得到“两个输入是否相同”。

\section{四、门网络构成组合逻辑}

组合逻辑的特点很直接：输出只由当前输入决定，不保存过去。它像一个硬件函数，但这个函数不是一步一步执行的过程，而是一张由门和连线组成的网络。输入变化后，信号沿多条路径同时传播；等延迟过去，输出稳定到当前输入对应的值。

先做一个三传感器判断器。输入 A、B、C 表示三个传感器，要求至少两个传感器为 1 时输出 1。这个需求不能只凭直觉画线，先列出行为：

\begin{longtable}{p{0.1\linewidth}p{0.1\linewidth}p{0.1\linewidth}p{0.16\linewidth}p{0.44\linewidth}}
\toprule
A & B & C & Y & 原因 \\
\midrule
0 & 0 & 0 & 0 & 一个条件都没有 \\
1 & 0 & 0 & 0 & 只有一个条件 \\
0 & 1 & 0 & 0 & 只有一个条件 \\
0 & 0 & 1 & 0 & 只有一个条件 \\
1 & 1 & 0 & 1 & A 和 B 同时满足 \\
1 & 0 & 1 & 1 & A 和 C 同时满足 \\
0 & 1 & 1 & 1 & B 和 C 同时满足 \\
1 & 1 & 1 & 1 & 三个都满足 \\
\bottomrule
\end{longtable}

从表中可以直接得到一个门网络：

\begin{lstlisting}
Y = (A AND B) OR (A AND C) OR (B AND C)
\end{lstlisting}

这条表达式里有三条“证据路径”。第一条路径证明 A、B 同时为 1；第二条路径证明 A、C 同时为 1；第三条路径证明 B、C 同时为 1。最后的 OR 不是任意合并，而是在检查三条证据中是否至少有一条成立。按这种方式阅读，门网络就不只是公式变形，而是把真值表中每一种输出为 1 的原因接成硬件。

两个 bit 相加是第二个基本组合逻辑例子。输出不止一个 bit，因为 1+1 得到二进制 10，需要一个和位 S 和一个进位 C。

\begin{longtable}{p{0.16\linewidth}p{0.16\linewidth}p{0.22\linewidth}p{0.22\linewidth}}
\toprule
A & B & S & C \\
\midrule
0 & 0 & 0 & 0 \\
0 & 1 & 1 & 0 \\
1 & 0 & 1 & 0 \\
1 & 1 & 0 & 1 \\
\bottomrule
\end{longtable}

观察表格可以看到，S 在 A、B 不同时为 1，C 在 A、B 同时为 1：

\begin{lstlisting}
S = A XOR B
C = A AND B
\end{lstlisting}

这就是半加器。它已经完成了一位相加，但它还缺一个输入：来自低位的进位。只要要做多位加法，最低位之外的每一位都不能只看 A 和 B，还要看低一位是否进上来。

全加器输入为 A、B、Cin，输出为 S、Cout。可以先把 A 和 B 用半加器相加，得到临时和 S1 和进位 C1；再把 S1 和 Cin 用第二个半加器相加，得到最终和 S 和第二个进位 C2；最后把两个进位合并。

\begin{lstlisting}
S1 = A XOR B
C1 = A AND B
S  = S1 XOR Cin
C2 = S1 AND Cin
Cout = C1 OR C2
\end{lstlisting}

把全加器一位一位串起来，就得到串行进位加法器。它的逻辑很清楚：第 0 位算出进位，传给第 1 位；第 1 位再把进位传给第 2 位。问题也很清楚：若最低位产生的进位一路传到最高位，最高位必须等前面所有进位稳定后才能稳定。这是组合逻辑里第一次真正遇到“逻辑正确”和“速度足够”的差别。真值表保证最后结果正确，传播延迟决定结果多久之后才可靠。

多路选择器解决“多个输入选一个输出”。2 选 1 选择器的行为是：

\begin{lstlisting}
if sel = 0, Y = A
if sel = 1, Y = B
\end{lstlisting}

上述 \code{if} 只是描述行为，电路本身没有顺序执行。它的门级表达式是：

\begin{lstlisting}
Y = (NOT sel AND A) OR (sel AND B)
\end{lstlisting}

选择器的作用很基础。若一条输出线可能来自加法器，也可能来自传感器，也可能来自某个寄存器，那么必须有一个明确的选择信号决定哪一路接到输出。没有选择器，多路来源同时驱动同一条线就会冲突。

译码器做相反方向的事：把一个二进制编号展开成 one-hot 信号。2-bit 输入可以选中 4 条输出中的一条。

\begin{longtable}{p{0.18\linewidth}p{0.48\linewidth}}
\toprule
输入 & 输出 \\
\midrule
00 & Y0=1，Y1/Y2/Y3=0 \\
01 & Y1=1，其他为 0 \\
10 & Y2=1，其他为 0 \\
11 & Y3=1，其他为 0 \\
\bottomrule
\end{longtable}

选择器把多路收成一路，译码器把编号展开成多路选择。后面无论是寄存器阵列、数据通路，还是存储器地址选择，都会反复看到这两个结构。

\section{五、延迟、负载、毛刺与验收}

门电路在纸上像瞬间给出输出的函数，真实硬件却至少有四个限制。第一，输入变化后，输出需要传播延迟才稳定。第二，一个输出能驱动的后级输入数量有限，后级太多会让切换变慢，甚至让电平余量变差。第三，门的输出必须重新形成清楚的 0/1，不能只靠一个偏弱路径勉强改变节点电压。第四，多条路径延迟不同，输入变化过程中可能出现短暂毛刺。

\begin{longtable}{p{0.2\linewidth}p{0.38\linewidth}p{0.32\linewidth}}
\toprule
限制 & 含义 & 后续影响 \\
\midrule
传播延迟 & 输出不会在输入变化的同一瞬间稳定 & 组合逻辑层数越多，稳定越慢 \\
扇出 & 一个输出只能可靠驱动有限后级 & 大网络需要缓冲器和分层布线 \\
电平恢复 & 输出要重新形成清楚 0/1 & 不能只追求逻辑表达式正确 \\
毛刺 & 不同路径到达时间不同 & 后级不能在毛刺期间采样 \\
\bottomrule
\end{longtable}

考虑一个反例。若把十几个门的输入全接到同一个弱输出上，真值表仍然可以保持正确，但物理电路可能切换缓慢。输出从 0 变 1 时，要给所有后级输入电容充电；如果下一级在它尚未进入可靠高电平区间时读取，就可能看到旧值或不确定值。因此，分析组合逻辑时不能只数门的个数，还要考察最长路径、负载和稳定时间。

短暂毛刺也很常见。假设一个输出由多条逻辑路径合成，输入同时变化时，每条路径稳定的时间可能不同。真值表只告诉我们变化前和变化后的稳定结果，却不保证中间不会短暂出现错误电平。如果后级只是另一个组合门，毛刺可能很快消失；如果后级在毛刺期间被时钟采样，错误就可能被保存下来。后面讲触发器和时钟时，会把这个问题放大讨论。

还要区分反相器和缓冲器。反相器改变逻辑含义，缓冲器通常保持逻辑含义但增强驱动能力或隔离负载。很多时候，插入缓冲器不是为了“多一个门”，而是为了让一个信号能推得动更长的线或更多的后级。软件读者容易觉得保持值不变的缓冲器没有意义，但在硬件里，驱动能力本身就是功能的一部分。

到这里，本章完成了从电平到组合逻辑的闭环：节点要有明确驱动路径；电平要有噪声余量；二极管提供方向性但不能恢复电平；受控开关让信号控制电流路径；路径表和真值表共同定义门；门网络构成加法器、选择器、译码器等组合电路；组合电路虽然没有记忆，却有延迟、负载和毛刺。下一章开始讨论另一件事：怎样让硬件保存状态，并按时钟边界一步一步推进。

\begin{classicnote}
本章的验收标准不是“会背 AND、OR、NOT、XOR 的真值表”，而是能把每个抽象还原成机器证据。

\begin{itemize}
\item 电平证据：前一级的最差高/低输出，必须落在后一级输入阈值允许的范围内，并留下噪声余量。
\item 路径证据：任何输出节点都要能写成“高电平路径是否导通、低电平路径是否导通、是否存在冲突电流、是否存在悬空”四项。
\item 门级证据：任选一个二输入门，列出 4 行真值表，再为每一行写出上拉网络和下拉网络状态。
\item 组合证据：半加器的 Sum 是 XOR，Carry 是 AND；全加器是在同一张真值表中纳入进位输入后的组合电路。
\item 延迟证据：同一个逻辑函数可以有不同门级网络，最长路径不同，稳定时间也不同。
\item 负载证据：一个门的输出不是无限共享变量，后级数量、线长和缓冲策略会改变电气行为。
\item 检查题：若组合网络输出短暂从 0 跳到 1 又回到 0，而最终值正确，它是否一定构成设计错误？答案取决于后级是否在毛刺期间采样。
\end{itemize}

经典系统教材常把抽象位函数和实际机器证据放在同一页上看。只看位函数会漏掉电平、驱动和时序；只看器件又会失去逻辑行为。两者合在一起，才是硬件设计对象。
\end{classicnote}
